\SourceLeaf{195}before the public. Perhaps, when I gain
more leisure than my present circumstances allow, I shall myself
make my investigations known in the future.

With respect I remain
\begin{center}
\begin{tabular}{@{}l@{\qquad}r@{}}
 & Your Honour's\\
Göttingen, 8 November & most obedient servant\\
\multicolumn{1}{c}{1824.} & \textsc{C. F. Gauss.}
\end{tabular}
\end{center}
\Needspace{9\baselineskip}
\medskip

\begin{center}
\textsc{Gauss} to \textsc{Olbers}. Göttingen, 3 May 1827.
\end{center}

.\ .\ .\ .\ .\ .\ .\ About 6 weeks ago I had the pleasure of meeting Mr \textsc{Dirichlet}
in person here. I believe I have written to you about him once before.
I mentioned to him the poor essay by
\textsc{Ivory}. He did not know it himself, but told me that Mr \textsc{Fourier}
had told him it was »nonsense« and that he (F.) had laughed heartily at a
fawning article about it in Férussac. He also judged his
memoir on the equilibrium figure of a rotating homogeneous
fluid much as I did. I was surprised to find, in
several other respects concerning people and matters, an extraordinary
agreement between Mr \textsc{Fourier}'s judgement and mine, for\,example
about imaginary quantities and about the impossibility of proving geometry a priori, and so
on. When I had briefly indicated to Mr \textsc{Dirichlet}, for\,example, my view on the latter,
he said that \textsc{Fourier} had expressed his own view to him in almost the
same words.
.\ .\ .\ .\ .\ .

